\chapter{Conclusion}\label{chap:conclusion}

\section{Summary of Achievements}

A charging station with one week of records offers few examples from which to learn its daily demand pattern. This thesis examined how much a forecaster can learn from established locations and whether recent forecast errors can tell it when to relax the influence of that prior training.

Source initialization gave the clearest improvement. On twelve UrbanEV targets, a GRU initialized from other locations reduced mean absolute error by 33.35\%, 31.64\% and 29.80\% with 7, 14 and 28 days of target history, respectively. Source initialization improved the seed-averaged error at all twelve targets for the seven- and fourteen-day budgets and at eleven targets for the twenty-eight-day budget. Both the transferred GRU and the target-only GRU had their initial training budgets selected on separate calibration targets. Even with 28 days of local observations, the transferred model retained an advantage for next-day hourly forecasts.

Once target training begins, the retention coefficient $\lambda$ controls how strongly the target encoder is penalized for departing from the source representation. On four calibration targets, $\lambda \approx 0.798$ gave the lowest error among 21 tested values. In the historical evaluation, switching coefficients on the basis of recent errors and their uncertainty improved on a constant $\lambda=1$ in nine of 21 dataset--city and history combinations. The rule improved on Online-best in eleven of the same 21 combinations. On twelve additional UrbanEV targets, a confidence rule selecting among the same candidate trajectories as the calibrated fixed control achieved lower error at 28 days and higher error at 7 and 14 days. These comparisons identify settings in which switching helped. They leave open how to recognize those settings before the next forecast.

A second approach combined forecasts from two GRUs with the same learning rate but different retention coefficients and online training windows. Both static and adaptive weighting improved on a single fixed-retention GRU in 18 of the 21 historical groups. With the same two expert forecasts, adaptive fusion improved on the calibrated static blend in twelve groups. On twelve UrbanEV Development targets at fourteen days, equal weighting reproduced the adaptive MAE within 0.000007 kWh. Historical comparisons with equal weighting favored the adaptive rule in all 21 groups, with reductions of at most 0.263\%. A follow-up evaluation fixed the methods and selected 46 additional Development identifiers before examining their outcomes. At 14 days of history, adaptive weighting reduced error by 0.67\% against the single fixed model and 0.36\% against the static mixture, giving each dataset--city group equal weight. Both 97.5\% uncertainty intervals included zero. The follow-up therefore leaves the incremental benefit of adapting the weights uncertain.

\section{Limitations and Future Work}

We simulated cold start by restricting how much recorded history the forecaster could use. A station's first day in an evaluation window may fall well after it opened. CHARGED also includes values interpolated or filled by the data publisher \cite{charged2025}. Our replay enforces label availability on the released series; the real-time availability of each upstream estimate cannot be reconstructed from that series. Following stations with verified opening dates would test whether source initialization also helps during the first weeks of operation, when customers are still discovering the site.

The historical groups were revisited during method development. The 46 additional identifiers provided a separate evaluation at the same seven dataset--city groups and during the same calendar period. Johannesburg and S\~ao Paulo each contributed a single identifier to that follow-up, which limits what their results can tell us about variation within either city. Evaluating a prespecified rule in a later period or in an additional city would test how well the observed effects carry over.

The one-step follow-versus-hold diagnostic found a mean advantage of only 0.000182 kWh on the calibration targets. A prospective test on additional targets or a later period would record each proposed change in $\lambda$ before the next block of forecasts. Starting from the same model state, one copy would adopt the proposed coefficient and another would keep the current value. Scoring both on the same subsequent days would show whether the evidence available at decision time predicts a benefit from switching. That prospective comparison would address the question left open here: when can a station with little history safely rely less on its source model?
